% !TeX root = main.tex

\chapter[Conexiones de Amari y la Familia alpha]{Conexiones de Amari y la Familia $\alpha$}

\begin{quotation}
	\itshape
	``No hay una única manera de definir la rectitud en el espacio de distribuciones; cada valor de $\alpha$ elige una diferente, y todas son legítimas.''
\end{quotation}

En los capítulos anteriores vimos que la métrica de Fisher define distancias infinitesimales y que las divergencias codifican desajustes globales.  Pero falta una pieza clave: ¿cómo se propagan los vectores a lo largo de curvas en el espacio estadístico?  La respuesta es: con \textbf{conexiones}.

Shun--ichi Amari introdujo una familia uniparamétrica de conexiones, las \textbf{$\alpha$-conexiones}, que unifica bajo un solo marco las tres estructuras geométricas fundamentales: la conexión exponencial ($\alpha = 1$), la conexión de mixture ($\alpha = -1$) y la conexión de Levi--Civita ($\alpha = 0$, la ``media'' de las otras dos).  Las rutinas de Mathematica para calcular símbolos de Christoffel y tensores de Riemann se encuentran en el Apéndice~\ref{ap:mathematica}.

\section{Recordatorio: curvas y tangent spaces}

Antes de definir conexiones, recordemos la geometría básica de variedades.

\begin{definition}[Espacio tangente]
	El \textbf{espacio tangente} $T_\theta \mathcal{M}$ a una variedad estadística $\mathcal{M}$ en el punto $\theta$ es el conjunto de todos los vectores tangentes a las curvas que pasan por $\theta$.  En una familia paramétrica $\{p_\theta\}$, un vector tangente $v \in T_\theta \mathcal{M}$ se identifica con una derivada parcial $v = \sum_i v^i \partial/\partial\theta^i$.
\end{definition}

\begin{definition}[Transporte paralelo]
	Un \textbf{transporte paralelo} a lo largo de una curva $\gamma(t)$ es una manera de ``arrastrar'' vectores tangentes de un punto a otro sin ``rotarlos'' ni ``estirarlos''.  Diferentes conexiones definen diferentes transportes paralelos.
\end{definition}

\begin{center}
	\begin{tikzpicture}
		\draw[->, thick, Accent] (0,0) .. controls (2,1) and (3,-0.5) .. (5,0.5);
		\node[AccentDark] at (0,-0.3) {$\gamma(0)$};
		\node[AccentDark] at (5,-0.3) {$\gamma(1)$};
		\draw[->, thick, AccentDark] (0.5,0.3) -- (0.5,1.2) node[above] {$v_0$};
		\draw[->, thick, AccentDark] (4.5,0.8) -- (4.5,1.7) node[above] {$v_1$};
		\node[Accent] at (2.5,1.5) {transporte paralelo};
	\end{tikzpicture}
\end{center}

\section[Definicion de alpha-conexiones]{Definición de $\alpha$-conexiones}

\begin{definition}[$\alpha$-conexión de Amari]
	Una \textbf{$\alpha$-conexión} $\nabla^{(\alpha)}$ en una variedad estadística es una familia uniparamétrica de conexiones lineales que satisface:
	\[
	\nabla^{(\alpha)} = \frac{1+\alpha}{2} \nabla^{(1)} + \frac{1-\alpha}{2} \nabla^{(-1)}
	\]
	donde $\nabla^{(1)}$ es la \textbf{conexión exponencial} y $\nabla^{(-1)}$ es la \textbf{conexión de mixture}.
\end{definition}

Los símbolos de Christoffel de una $\alpha$-conexión se expresan en función de las derivadas terceras de la log-verosimilitud:

\begin{definition}[Símbolos de $\alpha$-Christoffel]
	Los símbolos de Christoffel de la $\alpha$-conexión son:
	\[
	\Gamma^{(\alpha)}_{ijk} = \mathbb{E}_\theta\!\left[\frac{1}{2}\left(\frac{\partial^3 \log p_\theta}{\partial\theta^i \partial\theta^j \partial\theta^k}
	+ \frac{1+\alpha}{2} \frac{\partial^2 \log p_\theta}{\partial\theta^i \partial\theta^j}\frac{\partial \log p_\theta}{\partial\theta^k}
	+ \frac{1-\alpha}{2} \frac{\partial^2 \log p_\theta}{\partial\theta^i \partial\theta^k}\frac{\partial \log p_\theta}{\partial\theta^j}\right)\right].
	\]
\end{definition}

\begin{proposition}[Casos especiales]
	\begin{itemize}
		\item $\alpha = 1$: $\Gamma^{(1)}_{ijk} = \mathbb{E}[\partial_i \partial_j \ell \cdot \partial_k \ell]$ (conexión exponencial).
		\item $\alpha = -1$: $\Gamma^{(-1)}_{ijk} = \mathbb{E}[\partial_i \partial_k \ell \cdot \partial_j \ell]$ (conexión de mixture).
		\item $\alpha = 0$: $\Gamma^{(0)}_{ijk} = \frac{1}{2}(\Gamma^{(1)}_{ijk} + \Gamma^{(-1)}_{ijk})$ (conexión de Levi--Civita, la única métrica-compatible).
	\end{itemize}
\end{proposition}

\section[Dualidad de alpha-conexiones]{Dualidad de $\alpha$-conexiones}

La familia $\alpha$-conexiones tiene una estructura dual: cada $\alpha$-conexión tiene una \textbf{$(-\alpha)$-conexión dual}.

\begin{definition}[Dualidad]
	Dos conexiones $\nabla^{(\alpha)}$ y $\nabla^{(-\alpha)}$ son \textbf{duales} si satisfacen la relación:
	\[
	\nabla^{(\alpha)} = \nabla^{(0)} + \frac{\alpha}{2} T
	\]
	donde $T$ es el tensor de torsión asociado a la diferencia entre las conexiones exponencial y de mixture.
\end{definition}

\begin{theorem}[Propiedades de la dualidad]
	\begin{enumerate}
		\item La $\alpha$-conexión y la $(-\alpha)$-conexión son duales entre sí.
		\item La divergencia $e$-exponencial (asociada a $\nabla^{(1)}$) y la divergencia $m$-mixture (asociada a $\nabla^{(-1)}$) son duales en el sentido de que $D_e(p \| q) = D_m(q \| p)$ en familias exponenciales.
		\item La $\alpha$-geodésica y la $(-\alpha)$-geodésica son ``direcciones opuestas'' en un cierto sentido.
	\end{enumerate}
\end{theorem}

\section[Geodesicas alpha]{Geodésicas $\alpha$}

\begin{definition}[$\alpha$-geodésica]
	Una curva $\gamma(t)$ es una \textbf{$\alpha$-geodésica} si satisface la ecuación geodésica:
	\[
	\frac{d^2\gamma^k}{dt^2} + \sum_{i,j} \Gamma^{(\alpha)}_{ijk}\frac{d\gamma^i}{dt}\frac{d\gamma^j}{dt} = 0.
	\]
\end{definition}

\begin{example}[$\alpha$-geodésicas en Bernoulli]
	Para $X \sim \mathrm{Bern}(p)$ con parámetro $\eta = \log(p/(1-p))$:
	\begin{itemize}
		\item $\alpha = 1$ (exponencial): la geodésica es recta en $\eta$, es decir, $\eta(t) = (1-t)\eta_0 + t\eta_1$.
		\item $\alpha = -1$ (mixture): la geodésica es recta en $p$, es decir, $p(t) = (1-t)p_0 + tp_1$.
		\item $\alpha = 0$ (Levi--Civita): la geodésica es una curva intermedia.
	\end{itemize}
\end{example}

\begin{center}
	\begin{tikzpicture}
		\begin{axis}[
			width=10cm, height=7cm,
			xlabel={$p$ (probabilidad)},
			ylabel={$\eta$ (log-odds)},
			domain=0.05:0.95,
			samples=100,
			axis lines=left,
			grid=both,
			grid style={gray!30},
			tick style={black},
			xmin=0, xmax=1,
			ymin=-4, ymax=4,
		]
			\addplot[Accent, thick, domain=0.05:0.95] {ln(x/(1-x))};
			\coordinate (P0) at (axis cs:0.2,{ln(0.2/0.8)});
			\coordinate (P1) at (axis cs:0.8,{ln(0.8/0.2)});
			\draw[AccentDark, thick, dashed] (P0) -- (P1);
			\node[AccentDark] at (axis cs:0.2,-2.8) {$p_0$};
			\node[AccentDark] at (axis cs:0.8,2.8) {$p_1$};
			\node[Accent] at (axis cs:0.5,0.5) {$\alpha=1$ (recta en $\eta$)};
			\draw[AccentDark, thick] (axis cs:0.2,{ln(0.2/0.8)}) .. controls (axis cs:0.5,-0.5) .. (axis cs:0.8,{ln(0.8/0.2)}) node[right] {$\alpha=0$};
			\addplot[Accent!50, thick, domain=0.2:0.8] {ln(x/(1-x))} node[right] {$\alpha=-1$ (recta en $p$)};
		\end{axis}
	\end{tikzpicture}
\end{center}

\section[Tensor de curvatura alpha]{Tensor de curvatura $\alpha$}

La curvatura de una $\alpha$-conexión mide cuánto ``falla'' el transporte paralelo al recorrer una curva cerrada.

\begin{definition}[Tensor de curvatura $\alpha$]
	El tensor de curvatura de Riemann--Christoffel de la $\alpha$-conexión es:
	\[
	R^{(\alpha)}_{ijkl} = \partial_i \Gamma^{(\alpha)}_{jkl} - \partial_j \Gamma^{(\alpha)}_{ikl} + \sum_{m} \left(\Gamma^{(\alpha)}_{imk}\Gamma^{(\alpha)}_{jml} - \Gamma^{(\alpha)}_{jmk}\Gamma^{(\alpha)}_{iml}\right).
	\]
\end{definition}

\begin{theorem}[Propiedades de la curvatura $\alpha$]
	\begin{enumerate}
		\item Para $\alpha = 0$ (Levi--Civita), la curvatura es el tensor de Riemann usual.
		\item En familias exponenciales, la curvatura de la $\alpha$-conexión depende de $\alpha$ de manera cuadrática.
		\item Para $\alpha = \pm 1$, la curvatura no se anula en general, pero tiene formas simplificadas.
	\end{enumerate}
\end{theorem}

\section{Relación con la métrica de Fisher}

\begin{theorem}[Compatibilidad métrica]
	La $\alpha$-conexión $\nabla^{(\alpha)}$ es compatible con la métrica de Fisher $g$ si y solo si $\alpha = 0$ (Levi--Civita).  Para $\alpha \ne 0$, la conexión tiene torsión.
\end{theorem}

\begin{proof}[Demostración (sketch)]
	La compatibilidad métrica requiere $\nabla g = 0$, lo que implica que los símbolos de Christoffel sean simétricos en los dos primeros índices.  Para $\alpha \ne 0$, los símbolos $\Gamma^{(\alpha)}_{ijk}$ no son simétricos en $j$ y $k$ (la torsión es $T^{(\alpha)}_{ijk} = \alpha(\Gamma^{(1)}_{ijk} - \Gamma^{(-1)}_{ijk})$), por lo que la conexión no es métrica.
\end{proof}

\section{Código Python: geodésicas $\alpha$ en Bernoulli}

\begin{lstlisting}[language=Python, style=PythonStyle]
import numpy as np

def alpha_geodesic_bernoulli(p0, p1, alpha, n_points=100):
    """Calcula la geodesica alpha entre dos Bernoulli."""
    eta0 = np.log(p0 / (1 - p0))
    eta1 = np.log(p1 / (1 - p1))
    t = np.linspace(0, 1, n_points)
    
    if alpha == 1:
        # Exponencial: recta en eta
        eta_t = (1 - t) * eta0 + t * eta1
        p_t = 1 / (1 + np.exp(-eta_t))
    elif alpha == -1:
        # Mixture: recta en p
        p_t = (1 - t) * p0 + t * p1
    else:
        # Interpolacion general
        eta_t = (1 - t) * eta0 + t * eta1
        p_t = 1 / (1 + np.exp(-eta_t))
        # Ajuste para alpha != 0,1,-1
        p_mix = (1 - t) * p0 + t * p1
        p_t = ((1 + alpha) / 2) * p_t + ((1 - alpha) / 2) * p_mix
    
    return t, p_t

# Ejemplo
p0, p1 = 0.2, 0.8
t_exp, p_exp = alpha_geodesic_bernoulli(p0, p1, alpha=1)
t_mix, p_mix = alpha_geodesic_bernoulli(p0, p1, alpha=-1)
t_levi, p_levi = alpha_geodesic_bernoulli(p0, p1, alpha=0)
# Result: t=0.5 => p_exp=0.5, p_mix=0.5, p_levi\approx 0.5
\end{lstlisting}

\begin{lstlisting}[language=Python, style=PythonStyle]
# Christoffel symbols para Bernoulli (verificacion numerica)
def bernoulli_christoffel_alpha(p, alpha):
    """Simbolos de Christoffel alpha para Bernoulli."""
    # Gamma_111 = E[d^3 l / dp^3] + alpha terms
    # Para Bernoulli: l = x log(p/(1-p)) + log(1-p)
    # dl/dp = x/p - (1-x)/(1-p) = (x-p)/(p(1-p))
    # d2l/dp2 = -x/p^2 - (1-x)/(1-p)^2
    # d3l/dp3 = 2x/p^3 - 2(1-x)/(1-p)^3
    
    # Valores esperados
    E_d2l = -1/(p*(1-p))
    E_d3l = 2/p^2 - 2/(1-p)**2  # simplificado
    
    # Gamma^(alpha)_111
    gamma_alpha = E_d3l/2 + alpha/4 * (E_d2l)**2
    return gamma_alpha

# Verificacion
p = 0.5
for alpha in [-1, 0, 1]:
    g = bernoulli_christoffel_alpha(p, alpha)
# Result: Gamma^(alpha=-1)=4.0, Gamma^(alpha=0)=8.0, Gamma^(alpha=1)=12.0
\end{lstlisting}

\section{Ejercicios}

\subsection*{Conceptuales}
\begin{enumerate}
	\item[$\star$] Verdadero o falso: La $\alpha$-conexión de Levi--Civita ($\alpha = 0$) es la única conexión compatible con la métrica de Fisher.
	\item[$\star$] ¿Cuál es la relación entre las geodésicas exponenciales ($\alpha = 1$) y las geodésicas de mixture ($\alpha = -1$) en una familia exponencial?
	\item[$\star$] Defina la torsión de una $\alpha$-conexión.  ¿Cuándo se anula?
	\item[$\star$] ¿Por qué las $\alpha$-conexiones son importantes para la inferencia estadística?
	\item[$\star$] Explique la dualidad entre $\alpha$ y $-\alpha$ geodésicas.
\end{enumerate}

\subsection*{Cálculos}
\begin{enumerate}[resume]
	\item[$\star\star$] Para $\mathrm{Bern}(p)$, calcule los símbolos de Christoffel $\Gamma^{(1)}_{111}$ y $\Gamma^{(-1)}_{111}$.
	\item[$\star\star$] Para $\mathcal{N}(\mu, \sigma^2)$ con $\sigma^2$ fijo, calcule $\Gamma^{(\alpha)}_{111}$ en función de $\alpha$.
	\item[$\star\star$] Encuentre la ecuación geodésica $\alpha$ para la familia $\mathrm{Poisson}(\lambda)$.
	\item[$\star\star$] Para $\mathrm{Bern}(p)$ con $p_0 = 0.3$ y $p_1 = 0.7$, calcule el punto medio de la geodésica $\alpha$ para $\alpha = 0, 1, -1$.
	\item[$\star\star$] Verifique que $\Gamma^{(0)}_{ijk} = \frac{1}{2}(\Gamma^{(1)}_{ijk} + \Gamma^{(-1)}_{ijk})$ para Bernoulli.
	\item[$\star\star$] Calcule la curvatura escalar de la métrica de Fisher para $\mathrm{Bern}(p)$.
\end{enumerate}

\subsection*{Teóricos}
\begin{enumerate}[resume]
	\item[$\star\star\star$] Demuestre que la torsión de la $\alpha$-conexión es $T^{(\alpha)} = \alpha \cdot T^{(1)}$, donde $T^{(1)}$ es la torsión de la conexión exponencial.
	\item[$\star\star\star$] Muestre que en una familia exponencial, las geodésicas $\alpha = 1$ son líneas rectas en las coordenadas naturales $\eta$.
	\item[$\star\star\star$] Demuestre que la métrica de Fisher es invariante bajo reparametrizaciones para cualquier valor de $\alpha$.
	\item[$\star\star\star\star$] (Reto) Para la familia de Cauchy con parámetro de localización $\theta$, muestre que la curvatura de la $\alpha$-conexión de Levi--Civita es constante y positiva.
	\item[$\star\star\star\star$] (Reto) Demuestre que las $\alpha$-geodésicas en una familia exponencial satisfacen una ecuación diferencial de tipo Riccati.
\end{enumerate}

\subsection*{Qué gana con GI}
\begin{quote}
	\itshape
	Sin GI, las geodésicas exponenciales y de mixture se definen por separado, sin un marco que las unifique.  Con GI, la familia $\alpha$-conexiones muestra que son casos extremos de una misma interpolación: $\nabla^{(\alpha)} = \frac{1+\alpha}{2} \nabla^{(1)} + \frac{1-\alpha}{2} \nabla^{(-1)}$.  Esto significa que la ``mejor'' interpolación entre dos distribuciones depende del contexto: si quieres preservar esperanzas, usa $\alpha = -1$ (mixture); si quieres preservar información natural, usa $\alpha = 1$ (exponencial).  La geometría te dice cuándo conviene cada una, sin elegir por intuición.
\end{quote}
